% =====================================================================
%  MATH 55 -- Shared preamble for the Second Long Exam reviewers
%  Reproduces the look of "Lecture Notes 01" (Lecture Notes & Reviewer
%  No. 5): Latin Modern, small-caps running heads, ruled chapter titles,
%  amsthm environments numbered within chapters, framed strategy boxes.
%
%  Each document defines, BEFORE \input-ing this file:
%     \NotesNumber      e.g. 6
%     \NotesKind        e.g. Lecture Notes \& Reviewer No.~6
%     \NotesShortTitle  e.g. Directional Derivatives and Gradient
%     \NotesPdfTitle    plain-text title for the PDF metadata
% =====================================================================

% ---------- Fonts and encoding --------------------------------------
\usepackage[T1]{fontenc}
\usepackage[utf8]{inputenc}
\usepackage{lmodern}
\usepackage{microtype}

% ---------- Page geometry -------------------------------------------
\usepackage[a4paper,top=1in,bottom=1in,left=1.05in,right=1.05in,
            headheight=15pt,footskip=30pt]{geometry}
\setlength{\parindent}{0pt}
\setlength{\parskip}{6pt plus 2pt minus 1pt}

% ---------- Mathematics ---------------------------------------------
\usepackage{amsmath,amssymb,amsthm,mathtools}
\usepackage{bm}

% ---------- Graphics, tables, lists ---------------------------------
\usepackage{graphicx}
\usepackage[dvipsnames]{xcolor}
\usepackage{tikz}
\usetikzlibrary{arrows.meta,calc,angles,quotes,positioning}
\usepackage{pgfplots}
\pgfplotsset{compat=1.17}
\usepackage{booktabs,array,tabularx,multirow}
\usepackage{caption}
\usepackage{subcaption}
\captionsetup{font=small,labelfont=bf,width=0.92\linewidth}
\usepackage{enumitem}
\setlist{itemsep=2pt,topsep=4pt}
\usepackage{float}
\usepackage[most]{tcolorbox}

% ---------- Headings --------------------------------------------------
\usepackage{titlesec}
% "Chapter N" over a ruled small-caps title, each on a new page
\titleformat{\section}[display]
  {\normalfont\centering}
  {\large\scshape Chapter \thesection}
  {8pt}
  {\titlerule\vspace{10pt}\LARGE\scshape}
  [\vspace{4pt}\titlerule]
\titlespacing*{\section}{0pt}{0pt}{22pt}
\newcommand{\sectionbreak}{\clearpage}
\titleformat{\subsection}{\normalfont\large\bfseries}{\thesubsection}{1em}{}
\titleformat{\subsubsection}{\normalfont\normalsize\bfseries}{\thesubsubsection}{1em}{}
\setcounter{tocdepth}{2}

% Appendices are titled "Appendix A" instead of "Chapter A"
\newcommand{\startappendix}{%
  \appendix
  \titleformat{\section}[display]
    {\normalfont\centering}
    {\large\scshape Appendix \thesection}
    {8pt}
    {\titlerule\vspace{10pt}\LARGE\scshape}
    [\vspace{4pt}\titlerule]%
}

% ---------- Running heads and feet ----------------------------------
\usepackage{fancyhdr}
\newcommand{\NotesAuthor}{Juevesano, D.}
\fancypagestyle{m55}{%
  \fancyhf{}%
  \fancyhead[L]{\small\textsc{\NotesKind}}%
  \fancyhead[R]{\small\textsc{Math 55 -- \NotesShortTitle}}%
  \fancyfoot[L]{\small\NotesAuthor\ 2026 ALL RIGHTS RESERVED. Not for redistribution.}%
  \fancyfoot[R]{\small\thepage}%
  \renewcommand{\headrulewidth}{0.4pt}%
  \renewcommand{\footrulewidth}{0.4pt}%
}
% chapter-opening pages use "plain": keep the same head/foot there too
\fancypagestyle{plain}{%
  \fancyhf{}%
  \fancyhead[L]{\small\textsc{\NotesKind}}%
  \fancyhead[R]{\small\textsc{Math 55 -- \NotesShortTitle}}%
  \fancyfoot[L]{\small\NotesAuthor\ 2026 ALL RIGHTS RESERVED. Not for redistribution.}%
  \fancyfoot[R]{\small\thepage}%
  \renewcommand{\headrulewidth}{0.4pt}%
  \renewcommand{\footrulewidth}{0.4pt}%
}
\pagestyle{m55}

% ---------- Theorem-like environments --------------------------------
\theoremstyle{definition}
\newtheorem{definition}{Definition}[section]
\newtheorem{example}[definition]{Example}
\newtheorem{problem}[definition]{Problem}
\newtheorem{exercise}[definition]{Exercise}
\theoremstyle{plain}
\newtheorem{theorem}[definition]{Theorem}
\newtheorem{proposition}[definition]{Proposition}
\newtheorem{corollary}[definition]{Corollary}
\newtheorem{lemma}[definition]{Lemma}
\theoremstyle{remark}
\newtheorem{remark}[definition]{Remark}
\newtheorem{warning}[definition]{Warning}

\newenvironment{solution}{\begin{proof}[Solution]}{\end{proof}}
% bold step labels inside solutions, e.g. \step{1}{Partial derivatives.}
\newcommand{\step}[2]{\par\smallskip\noindent\textbf{Step #1: #2}\ }
\newcommand{\stepx}[1]{\par\smallskip\noindent\textbf{#1}\ }
\newcommand{\checkok}{\ \checkmark}

% ---------- Boxes -----------------------------------------------------
% Framed "Strategy" box (as on p. 7 of Lecture Notes 01)
\newtcolorbox{strategy}[1]{%
  enhanced,breakable,colback=white,colframe=black!75,boxrule=0.6pt,
  sharp corners,left=10pt,right=10pt,top=8pt,bottom=8pt,
  before upper={\textbf{#1}\par\smallskip}}
% Highlighted key-result box
\newtcolorbox{keybox}[1]{%
  enhanced,breakable,colback=black!3,colframe=black!75,boxrule=0.6pt,
  sharp corners,left=10pt,right=10pt,top=8pt,bottom=8pt,
  before upper={\textbf{#1}\par\smallskip}}

% "Alternative method" box: a quicker route shown AFTER the full solution
\newtcolorbox{altmethod}{%
  enhanced,breakable,colback=white,colframe=black!45,boxrule=0.5pt,
  sharp corners,left=10pt,right=10pt,top=6pt,bottom=6pt,
  before upper={\textbf{Alternative method (optional).}\ }}
% Boxed concluding sentence of a solution
\newcommand{\finalans}[1]{\par\smallskip\noindent\textbf{Final answer.}\ #1}

% ---------- Notation --------------------------------------------------
\newcommand{\R}{\mathbb{R}}
\newcommand{\vb}[1]{\mathbf{#1}}
\newcommand{\vi}{\mathbf{i}}
\newcommand{\vj}{\mathbf{j}}
\newcommand{\vk}{\mathbf{k}}
\newcommand{\vu}{\mathbf{u}}
\newcommand{\vv}{\mathbf{v}}
\newcommand{\vr}{\mathbf{r}}
\newcommand{\vx}{\mathbf{x}}
\newcommand{\vF}{\mathbf{F}}
\newcommand{\vT}{\mathbf{T}}
\newcommand{\vN}{\mathbf{N}}
\newcommand{\vB}{\mathbf{B}}
\newcommand{\vzero}{\mathbf{0}}
\newcommand{\ip}[1]{\left\langle #1 \right\rangle}
\DeclarePairedDelimiter{\norm}{\lVert}{\rVert}
\DeclarePairedDelimiter{\abs}{\lvert}{\rvert}
\newcommand{\dd}{\mathrm{d}}
\newcommand{\Du}{D_{\mathbf{u}}}
\newcommand{\grad}{\nabla}
\newcommand{\pd}[2]{\frac{\partial #1}{\partial #2}}
\newcommand{\vecPQ}{\overrightarrow{PQ}}
\newcommand{\eval}[1]{\Big|_{#1}}

% ---------- Interactive links (underlined, like Lecture Notes 01) ----
\newcommand{\extlink}[2]{\href{#1}{\underline{#2}}}
\newcommand{\interactive}[2]{%
  \par\noindent\textit{Interactive:} \extlink{#1}{#2}\ }
\newcommand{\desmosTD}{https://www.desmos.com/3d}
\newcommand{\desmosTwoD}{https://www.desmos.com/calculator}

% ---------- Hyperref (load last) -------------------------------------
\usepackage[hidelinks]{hyperref}
\hypersetup{pdfauthor={Juevesano, D.},pdftitle={\NotesPdfTitle}}

% ---------- Title page -----------------------------------------------
% \makenotestitle{Title}{Subtitle}{Line under "No. N"}{About text}
\newcommand{\makenotestitle}[4]{%
  \begingroup
    \thispagestyle{plain}
    \centering
    \vspace*{1.6cm}
    {\large\scshape University of the Philippines Cebu\par
     Mathematics Program\par}
    \vspace{1.6cm}
    \rule{\linewidth}{0.5pt}\par\vspace{0.8cm}
    {\Huge\scshape #1\par}
    \vspace{0.5cm}
    {\large #2\par}
    \vspace{0.8cm}\rule{\linewidth}{0.5pt}\par
    \vspace{1.4cm}
    {\large\scshape \NotesKind\par}
    \vspace{0.4cm}
    {#3\par}
    \vfill
    \begin{minipage}{0.84\linewidth}
      \small\raggedright #4
    \end{minipage}
    \vfill
    {\NotesAuthor\quad 2026\par}
    \vspace{0.6cm}
  \clearpage
  \endgroup
}

% The table of contents is a starred \section, so it automatically gets
% the centred, ruled small-caps title used for chapters.
